\section{Discussion}
Three observations explain why 20 parameters close this much of the gap. Counts supply the center, fitting corrects it, and the zero-center control isolates information from capacity. A 650-parameter Linear control uses over 32 times as many values but reaches higher NLL, while eleven-value BCTS with the same center approaches vector scaling with lower client ECE. The center, rather than the parameterization, carries the gain.

The upload half answers a different question: whether a contribution can be shown to have entered the model. A local table would store the same fields, but it would be written by the party that also decides which updates are aggregated, so the client's own record would prove little. Placing the binding on a ledger moves it outside the coordinator's local database write path, and an incorporation receipt issued by the coordinator ties the on-chain entry to the model state that authorizes the next residual. Retries then resume from a recorded state rather than from an assumption.

The accounting separates the costs that deployments trade differently. Personalization adds 80 B locally, one TRAIN forward pass and a twenty-dimen\-sional convex solve. Shared training requires 54,067 B uploads and a 547,496 B FP32 residual per client, with dense downlinks. Anchoring adds 341 B of bound context and authentication per upload and one registration per client--round, measured at 281,965 gas on average across fifteen registrations.
